% !TEX root = ../main.tex

\ustcsetup{
  keywords  = {大语言模型, 金融文本分析, 检索增强生成, 股票方向预测, 多源信息融合, 强化学习},
  keywords* = {Large Language Models, Financial Text Analysis, Retrieval-Augmented Generation, Stock Direction Prediction, Multi-source Information Fusion, Reinforcement Learning},
}

\begin{abstract}
投资者查看一只股票时，通常不会只看价格曲线。成交量和 K 线记录的是已经发生的交易，论坛讨论、新闻报道、机构研报和搜索结果中则会出现事件、情绪和预期变化。本文从这一使用习惯出发，设计并实现了 MarketMind。系统输入股票代码后，能够获取近期评论、新闻、研报、7 日 K 线和网页检索结果，并把这些材料整理成大语言模型可以阅读的分析输入，用于辅助判断后续走势。

系统分别处理 Comments、News 与 Report 三类文本，并在在线页面中接入 SerpAPI，补充公开网页中的最新线索。不同来源的内容先经过筛选和摘要，再交给模型综合判断。最终输出不再要求固定给出未来 1、3、7 日三个涨跌标签，而是给出股票未来趋势、支持上涨的因素、下行风险以及信息之间的分歧。由于历史时点的搜索结果很难完整复现，本文另行构造离线实验，比较外部文本加入前后、微调前后模型在未来 1 日、3 日和 7 日方向任务上的表现。实验以 Qwen2.5-7B-Instruct 为基础模型。

实验结果表明，只依赖历史 K 线时，模型表现接近随机二分类水平；加入新闻和研报摘要后，未来 1 日准确率由 45.49\% 提升到 52.38\%，未来 3 日准确率由 41.18\% 提升到 52.26\%，7 日预测变化不明显。经过 LoRA 监督微调后，Qwen2.5-7B-Instruct 在使用外部文本时取得最高 56.13\% 的预测成功率。这个结果不能说明模型已经具备稳定交易能力，但能看出文本材料主要在较短周期上改善了方向判断。本文还实现了面向未来 1、3、7 个交易日方向任务的 GRPO 训练入口，为后续继续优化模型输出保留了工程基础。
\end{abstract}

\begin{abstract*}
When investors look at an individual stock, they rarely rely on the price curve alone. K-line data and trading volume describe what has already happened, while forum comments, news reports, analyst reports, and search results often contain events, sentiment, and changing expectations. Starting from this use case, this thesis designs and implements MarketMind. Given a stock code, the system retrieves recent comments, news, reports, 7-day K-line data, and web search results, then turns them into analysis inputs that a large language model can read.

The system processes Comments, News, and Reports separately, and the online prototype integrates SerpAPI to add recent public web information. Content from each source is filtered and summarized before being passed to the model for a final judgment. Instead of forcing the online system to output three fixed labels for the next 1, 3, and 7 trading days, it now produces a future trend forecast, supporting factors, downside risks, and disagreements across sources. Since historical search results are difficult to reproduce exactly, this thesis uses a separate offline setting to compare model performance with and without external text, and before and after fine-tuning. Qwen2.5-7B-Instruct is used as the base model.

The offline results show that a model using only historical K-line data performs close to random binary classification. After news and report summaries are added, the 1-day accuracy increases from 45.49\% to 52.38\%, and the 3-day accuracy increases from 41.18\% to 52.26\%, while the 7-day result changes little. After LoRA supervised fine-tuning, Qwen2.5-7B-Instruct reaches a highest prediction success rate of 56.13\% when external text is used. This does not mean that the model is ready for trading, but it shows that text inputs mainly help shorter-horizon direction prediction in the current setting. This thesis also implements a GRPO training entry for 1-, 3-, and 7-trading-day direction prediction, leaving an engineering basis for further optimization.
\end{abstract*}
